\section{Ordinary launch, powered flight, and coast}
Run initialization selects active stages and motor installations, builds the initial event queue, establishes pad position and orientation, and initializes listeners. Orientation incorporates launch-rod inclination and bearing; the status also selects an initial roll orientation relative to the flight-condition convention. The initial-spin setting is a separate MIT addition.

The launch event can activate configured motors after their ignition delays. Ignition begins a motor's local clock; burnout ends thrust; ejection follows the selected delay where applicable. Liftoff and rod clearance are physical milestones distinct from ignition. Before liftoff the vehicle cannot translate downward through the pad, and while constrained by the rod the solver projects acceleration onto the guide direction.

\class{SimulationStatus} computes effective guide travel from configured rod length and the aftmost modeled launch-lug position. With no lugs it assumes a full-length tower launcher. The inspected loop specifically checks \class{LaunchLug}; the presence of a \class{RailButton} is not equivalent to that lug-length correction. The rod model constrains motion but does not simulate rail bending, button contact friction, or mechanical release loads.

After burnout, aerodynamic forces, gravity, wind, and changing configuration continue to govern coast. The engine detects apogee and altitude transitions, schedules recovery/staging rules, and records warnings and milestones. If deployment occurs before the natural coast apex, an auxiliary optimum-coast calculation can estimate the altitude and time that would have been reached without that deployment. This supports optimum-delay reporting; it is not a second observed flight history.

\section{What a numerical step means}
A Runge--Kutta step forms provisional positions, velocities, orientations, and rates at prescribed substep times, evaluates the derivatives there, and combines them into the accepted endpoint. The quaternion represents orientation without Euler-angle singularities, but it must remain normalized and used with the correct transform direction. The world angular-rate vector and body aerodynamic rates are related by rotation; they are not interchangeable arrays of three numbers.

The stepper records the force/mass/condition quantities associated with its accepted data point and supplies calculated events to the engine. Event-aware limits prevent stepping blindly beyond a known ignition, ejection, or other scheduled boundary. The retained RK4 and MIT RK6 implementations have different stage formulas and can have different limiting behavior; their names alone do not establish a formal adaptive local-error tolerance. The ordinary maximum timestep and angular-step settings express the actual available controls.

Exceptions separate cancellation, calculation failure, listener termination, and other run failures. System listeners such as ground-hit or recovery-stop listeners can serve special computations with different stopping rules. A headless caller must still interpret warnings/events and close run resources, rather than merely checking that the method returned an object.

\section{Tumble and grounded modes}
The tumble stepper is a translational empirical descent model. It combines symmetric-body projected area with fin planform area and fin-count efficiency factors. The coefficients in \class{BasicTumbleStepper} are 0.56 for the body contribution and 1.42 for the effective fin contribution, normalized to rocket reference area. The fin-count table approximates exposure during tumble. This is not rotational integration of a tumbling rigid body.

Landing and tumble use \class{AbstractEulerStepper}, while \class{GroundStepper} handles grounded continuation. The engine can change the stepper while retaining the branch's identity and event history. A held quaternion after deployment/tumble reflects this numerical model, not measured canopy alignment. These mode changes must be respected by plots, physical sensors, and extensions that continue through descent.

\sourceline{\source{core/src/main/java/info/openrocket/core/simulation/SimulationStatus.java}{initialization and guide length}; \source{core/src/main/java/info/openrocket/core/simulation/BasicEventSimulationEngine.java}{inherited event lifecycle}; \source{core/src/main/java/info/openrocket/core/simulation/BasicTumbleStepper.java}{BasicTumbleStepper}; \source{core/src/main/java/info/openrocket/core/simulation/GroundStepper.java}{GroundStepper}; \source{core/src/main/java/info/openrocket/core/simulation/listeners/system/OptimumCoastListener.java}{OptimumCoastListener}.}
